\documentclass[12pt]{article}

% a template that a friend gave, it's worked well enough for me
% i have added some packages and stuff that have proved useful

\usepackage{fancyhdr}
\usepackage{tipa}
\usepackage{fontspec}
\usepackage{amsfonts}
\usepackage{enumitem}
\usepackage[margin=1in]{geometry}
\usepackage{graphicx}
\usepackage{float}
\usepackage{amsmath}
\usepackage{braket}
\usepackage{amssymb}
\usepackage{booktabs}
\usepackage{hyperref}
\usepackage{mathtools}
\usepackage{xcolor}
\usepackage{float}
\usepackage{algpseudocodex}
\usepackage{titlesec}
\usepackage{bbm}

\pagestyle{fancy}
\fancyhf{} % sets both header and footer to nothing
\lhead{Kevin Sheng}
\setmainfont{Comic Neue}
\renewcommand{\headrulewidth}{1pt}
\setlength{\headheight}{0.75in}
\setlength{\oddsidemargin}{0in}
\setlength{\evensidemargin}{0in}
\setlength{\voffset}{-.5in}
\setlength{\headsep}{10pt}
\setlength{\textwidth}{6.5in}
\setlength{\headwidth}{6.5in}
\setlength{\textheight}{8in}
\renewcommand{\headrulewidth}{0.5pt}
\renewcommand{\footrulewidth}{0.3pt}
\setlength{\textwidth}{6.5in}
\usepackage{setspace}
\usepackage{multicol}
\usepackage{float}
\setlength{\columnsep}{1cm}
\setlength\parindent{24pt}
\usepackage [english]{babel}
\usepackage [autostyle, english = american]{csquotes}
\MakeOuterQuote{"}

\setlength{\parskip}{6pt}
\setlength{\parindent}{0pt}

\titlespacing\section{0pt}{12pt plus 4pt minus 2pt}{0pt plus 2pt minus 2pt}
\titlespacing\subsection{0pt}{12pt plus 4pt minus 2pt}{0pt plus 2pt minus 2pt}
\titlespacing\subsubsection{0pt}{12pt plus 4pt minus 2pt}{0pt plus 2pt minus 2pt}

\hypersetup{colorlinks=true, urlcolor=blue}

\newcommand{\correction}[1]{\textcolor{red}{#1}}

\usepackage{nicematrix}

\rhead{Math 110A}

\makeatletter
\def\@seccntformat#1{%
  \expandafter\ifx\csname c@#1\endcsname\c@section\else
  \csname the#1\endcsname\quad
  \fi}
\makeatother

\newcommand{\lra}{\xLeftrightarrow}
\newcommand{\ra}{\xRightarrow}
\DeclareMathOperator{\pr}{pr}
\DeclareMathOperator{\Mat}{Mat}

\begin{document}

\section{Problem 1}\label{sec:p1}

\subsection{Part A}\label{sec:p1a}

If we have an $\underline{f}$, then $\underline{f}(\pi(x))=f(x)$ for all $x \in X$.

Consider two functions $f_1$ and $f_2$ that both satisfy this property.
For any $x \in X$, we have that $f_1(\pi(x))=f(x)=f_2(\pi(x))$.
Since $\pi$ is surjective, $f_1(s)=f_2(s)$ for all $s \in X/R$.

\subsection{Part B}\label{sec:p1b}

I'll prove the backward direction and its converse.

Let $f$ be constant on $R$-equivalence classes.
Then, we can let $\underline{f}(\pi(x))=f(x)$ for all $x$.
This map is well-defined due to that $\pi$ is surjective and the following chain of logic:
\[\pi(x)=\pi(y) \implies f(x)=f(y) \implies \underline{f}(\pi(x))=\underline{f}(\pi(y))\]
Also by construction, $\underline{f}$ computes through the diagram given in the problem.

For the converse, suppose there existed $x_1, x_2 \in X: \pi(x_1)=\pi(x_2) \land f(x_1) \ne f(x_2)$.
Let $p=\pi(x_1)=\pi(x_2)$ for convenience.

Then a hypothetical $\underline{f}$ would need to satisfy
$\underline{f}(p)=f(x_1)$ and $\underline{f}(p)=f(x_2)$ at the same time,
which is an impossibility since $f(x_1) \ne f(x_2)$. $\square$

\pagebreak

\section{Problem 2}\label{sec:p2}

\subsection{Part A}

Following the hint, we take $Y=Q$ and $\phi: Q \to Q$ defined by $\phi(q)=q$.
With this configuration, we have $q^*(\phi)=\phi \circ q=q$.

By the premise, $q^*(\phi)$ must be constant on $R$-equivalence classes,
so $q$ must be as well.

\subsection{Part B}

The existence and uniqueness of $\beta$ is implied by \ref{sec:p1}.
It is $\beta(\pi(x))=q(x)$.

For $\chi$ to exist, we need $q$ to be a surjection.

We prove that it is by showing the contrapositive.
If $\exists w \in Q: q(x) \ne w\ \forall x \in X$, then we can let $Y=Q$
and consider the image of the following two maps:
\begin{gather*}
  \phi_1(\tilde{q})=\tilde{q} \\
  \phi_2(\tilde{q})=\begin{cases}
    q(x) & \tilde{q}=w\text{ ($x$ is any $x \in X$)} \\
    q    & \text{otherwise}
  \end{cases}
\end{gather*}
$\phi_1 \ne \phi_2$ since their behavior differs on $w$.

Then, $\phi_1 \circ q=q$, and $\phi_2 \circ q=q$ as well since no $x$
will ever hit the first case in $\phi_2$.
Thus, $q$ not being a surjection violates the preconditions of its existence.

We also need that $\pi(x_1) \ne \pi(x_2) \implies q(x_1) \ne q(x_2)$.

BWOC let $\pi(x_1) \ne \pi(x_2)$ and $q(x_1)=q(x_2)$.
Let $q(x_1)=w$ for convenience.
Now let $Y=W \cup \{w'\}$, where $w'$ is just some dummy element that I made up.

By the premise, $\exists \phi: Q \to Y$ s.t.
\[(\phi \circ q)(x)=\begin{cases}
    w    & \pi(x)=\pi(x_1)  \\
    w'   & \pi(x)=\pi(x_2)  \\
    q(x) & \text{otherwise}
  \end{cases}\]
By construction, this target function is constant on equivalence classes.
However, since $q(x_1)=q(x_2)$, we have $(\phi \circ q)(x_1)=(\phi \circ q)(x_2)$
no matter which $\phi$ we take.
Thus, no such $\phi$ that satisfies our target can exist and we have a contradiction.

With these two preconditions proven, we can finally define $\chi(q(x))=\pi(x)$.
As $q$ is a surjection, this indeed defines a value for all elements of $Q$.
There's also no conflicts in any definitions, as $q(x_1)=q(x_2) \implies \pi(x_1)=\pi(x_2)$.

For uniqueness, note that for any $\chi_1$ and $\chi_2$ satisfying the given properties,
$\chi_1(q(x))$ and $\chi_2(q(x))$ both have the equal $\pi(x)$.
As $q$ is surjective, this equality must hold for all $w \in Q$ and therefore $\chi_1=\chi_2$.

\subsection{Part C}

It STP that $\text{id}=\chi \circ \beta=\beta \circ \chi$.

We have
\begin{align*}
  (\chi \circ \beta)(\pi(x))
   & = \chi(\beta(\pi(x))) \\
   & = \chi(q(x))          \\
   & = \pi(x)
\end{align*}
and
\begin{align*}
  (\beta \circ \chi)(q(x))
   & = \beta(\chi(q(x))) \\
   & = \beta(\pi(x))     \\
   & = q(x)\quad\square
\end{align*}

\pagebreak

\section{Problem 3}

\subsection{Part A}

The proof of this is the same as that in \ref{sec:p1a}.

\subsection{Part B}

The proof is almost the same as \ref{sec:p1b}, but it's different enough to warrant a writeup.

I'll prove the backward direction and its converse again.

If $f$ annihilates $S$, then $s \in S \implies f(s)=0$.

Let us define $\underline{f}(\pi(v))=f(v)$.
This map is well-defined since $\pi$ is a surjection onto $V/S$ and
\begin{align*}
  \pi(v_1)=\pi(v_2)
   & \implies v_1-v_2 \in S   \\
   & \implies f(v_1-v_2)=0    \\
   & \implies f(v_1)-f(v_2)=0 \\
   & \implies f(v_1)=f(v_2)
\end{align*}
By construction, this satisfies the commutativity diagram as well.

For the backwards direction, BWOC say $\exists s \in S: f(s) \ne 0$.

There's two cases after this.
If $S=V$, then $V/S$ would only have one element, the zero vector.
Since $\underline{f}$ is a linear map, $\underline{f}(\pi(v))=0$ for all $v$.
However, we also need $\underline{f}(\pi(s))=f(s)$, which isn't equal to $0$.

OTOH, if $S \subsetneq V$, then $\exists v \in V \setminus S$.
By definition, $\pi(v)=\pi(v + s)$, so $\underline{f}(\pi(v))=\underline{f}(\pi(v+s))$.
However, at the same time, $f(v+s)=f(v)+f(s) \ne f(v)$, since the two
being equal would imply $f(s)=0$.
This forces $\underline{f}(\pi(v)) \ne \underline{f}(\pi(v+s))$,
which still results in the desired contradiction. $\square$

\pagebreak

\section{Problem 4}

DEAR GOD WHY IS IT ALL FUNCTIONS

But yeah, we know there's a unique linear $\beta: V/S \to Q$ that commutes with the diagram
defined by $\beta(\pi(v))=f(v)$.

I'll prove the existence of an inverse $\chi: Q \to V/S$.

\subsection{\texorpdfstring{$q$}{q} is Surjective}

First we show that like in \ref{sec:p2}, $q$ must be surjective.
BWOC say $\exists w \notin q(V)$.
This then means that $w$ is linearly independent of any basis of $q(V)$.

Let $(v_i)_{i \in I}$ be a basis for $q(V)$, and $(w_j)_{j \in J}$
be whatever elements are needed to extend the basis to $Q$.
By 115 probably, we can force $w_1=w$.

Now consider these linear transformations from $Q$ to itself:
\begin{gather*}
  \phi_1(w_i)=\begin{cases}
    0   & i=1              \\
    w_i & \text{otherwise}
  \end{cases} \\
  \phi_1(v_i)=v_i \\
  \phi_2(w)=w
\end{gather*}
$\phi_1$ is a well-defined linear transformation since I've defined its actions on a basis of $Q$,
and $\phi_2$ is just the identity map.
These two linear maps are clearly distinct,
as they behave differently on $w_1=w$.

Now let us consider what happens when we compose them with $q$:
\begin{align*}
  (\phi_1 \circ q)(v)
  &= \phi_1\left(\sum_{j \in J} \lambda_j v_j\right) \\
  &= \sum_{j \in J} \lambda_j v_j \\
  &= \phi_2\left(\sum_{j \in J} \lambda_j v_j\right) \\
  &= (\phi_2 \circ q)(v)
\end{align*}
This means that $q^*(\phi_1)=q^*(\phi_2)$, a contradiction.

\pagebreak

\subsection{Injective Over \texorpdfstring{$V/S$}{V/S}}

We also need that $\pi(v_1) \ne \pi(v_2) \implies q(v_1) \ne q(v_2)$.

Again, BWOC say $\exists v_1, v_2: v_1 - v_2 \notin S \land q(v_1)=q(v_2)$.
This means that $v_1 - v_2 \in \ker(q)$ while at the same time
being linearly independent of any basis of $S$.
Let's call this rogue vector $v$.

Now let us construct a basis for $V$ defined by the union of $(s_i)_{i \in I}$,
a basis for $S$, and $(v_j)_{j \in J}$ with $v_1=v$, which contains
whatever vectors are left to extend $(s_i)_{i \in I}$ to a basis for $V$.

By our premise, if we let $W=\text{span}((v_j)_{j \in J})$,
there must be some linear map $\phi: Q \to W$ s.t.
\begin{gather*}
  (\phi \circ q)(s_i)=0 \\
  (\phi \circ q)(v_i)=v_i
\end{gather*}
By definition, this linear map annihilates $S$ since we send all of its
basis vectors to $0$.

However, no matter what we set $\phi$ to, we'll always have $\phi(q(v_1))=\phi(0)=0$,
thus making it impossible for $\phi(q(v_1))=v_1$.
Thus, $\pi(v_1) \ne \pi(v_2) \implies q(v_1) \ne q(v_2)$.

\subsection{Putting it Together}

Now we can finally define $\chi(q(v))=\pi(v)$.
By similar reasoning as in \ref{sec:p2}, this mapping indeed defines
a value for every $q \in Q$ and ensures that no input gets forced
to two different outputs.

We can also use basically the same chain of equalities previously
used to show that $\chi \circ \beta=\beta \circ \chi$, thus making
$V/S$ isomorphic to $Q$. $\square$

\pagebreak

\section{Problem 5}

BWOC let there be a countable basis for $\C^\infty$.

Since only finite linear combinations are allowed, let us consider
all finite subsets of this countable basis along with their spans.
Since we force subsets to be finite, there's only a countable amount of them:
$S_1, S_2, S_3, \cdots$.

Consider any subspace $S_i$ with dimension $k < \infty$
along with the following linear transformation $T_{s,l}(x_1, x_2, \cdots)$ with $s,l \in \N$:
\[T_{s,l}(x_1, x_2, x_3, \cdots) = \sum_{i=s}^{s+l} x_i\]
In actual English, $T_{s,l}$ is the projection from $\C^\infty$ onto the $s$th through $s+l$th coordinates.

The codomain of $T_{s,k}$ has dimension $k+1$, so there must exist some element
in there that isn't hit by the image of $S_i$ under $T_{s,k}$:
$\exists b \in T_{s,k}\left(C^\infty\right) \setminus T_{s,k}(S_i)$.

Then, by construction, for any $v \in C^\infty$ s.t. $v_{s-1+i}=b_i$, $v \notin S_i$.

Now we can diagonalize by using the $b$ for every subspace $S_i$ under some partitioned coordinates.
To do this, define the following things:
\begin{gather*}
  k_i=\dim S_i \\
  s_i=1+\sum_{j=1}^{i-1} (k_j+1) \\
  T'_i=T_{s_i,k_i}
\end{gather*}
These linear projections never use the same coordinate twice, and each subspace's
associated $T'_i$ has a codomain of dimension $k+1$.
As we proved earlier, each $T'_i$ must have some associated $b_i$ s.t.
for any $v \in C^\infty$ s.t. $v_{s_i-1+j}=(b_j)_i$, $v \notin S_i$.

Thus, the vector
\[\begin{bmatrix}
  b_1 \\
  b_2 \\
  b_3 \\
  \vdots
\end{bmatrix}\]
is not in any $S_i$, which contradicts that our basis was a spanning set. $\square$

\end{document}
